%% toc: P2 13 13.4,13.5,13.6,13.7
\section{Tableau soundness and completeness}
\label{sec:p02-w02-p02-014-tableau-metatheory}

A propositional tableau is a search for a valuation satisfying a finite set
of formulas.  The search either closes every branch, which certifies that no
such valuation exists, or ends with an open branch from which a satisfying
valuation can be read off.  This section makes that statement exact for the
calculus of \cref{def:p02-cpl-tableau}, whose rules are displayed on the
calculus sheet \texttt{sheet:\allowbreak tableaux-cpl}, and proves each of its parts: soundness of
closure, termination of exhaustive expansion, the truth lemma for open
saturated branches, and completeness.  Together these results establish every
clause of \cref{thm:p02-tableau-adequacy}.  It then shows why the
termination argument is special to the propositional case and treats the
first-order calculus of \cref{def:p02-fol-tableau}, extended by equality
rules: soundness, extraction of a Herbrand countermodel from an open
saturated (fair) branch, completeness, and a fair systematic procedure that
halts exactly on unsatisfiable inputs.  These results establish every clause
of \cref{thm:p02-fol-tableau-adequacy}.

\emph{Inputs.}  A finite set \(\Gamma\subseteq\mathsf{Fm}_{\mathsf{CPL}}\)
in the language of \cref{sec:prop-syntax-semantics-01}, with primitive
\(\bottom,\neg,\land,\lor,\to\); and, for consequence and validity, one
further formula \(\varphi\).  \emph{Conclusions.}  Either a closed tableau,
which shows that \(\Gamma\) is unsatisfiable (or that \(\Gamma\entails\varphi\)),
or a finished tableau with an open branch together with an explicit valuation
that satisfies \(\Gamma\) (or satisfies \(\Gamma\) and falsifies
\(\varphi\)).  Valuations, satisfaction \(v\entails\chi\), satisfiability,
validity and semantic consequence are those of
\cref{sec:prop-syntax-semantics-02}; \(\widehat v\) is the extension of a
valuation \(v\) to all formulas.  In the first-order part
(\cref{sec:p02-w02-p02-014-fol-calculus,sec:p02-w02-p02-014-fol-herbrand,sec:p02-w02-p02-014-fol-systematic})
the input is a finite set of first-order sentences, with or without
equality, and the conclusion is either a finite closed tableau, which shows
the set unsatisfiable, or an open saturated branch, possibly infinite,
together with a structure built from closed terms that satisfies the set.

\subsection{Input convention, expandable occurrences and finished tableaux}
\label{sec:p02-w02-p02-014-input}

\begin{hbconvention}{def:p02-w02-p02-014-input}{Root branch and the three tests}
Fix a finite set \(\Gamma\) and list it without repetition in a fixed
order as \(\gamma_1,\ldots,\gamma_n\), where \(n=0\) is allowed.  The \emph{root branch}
\(B_\Gamma\) is the list \(\langle\gamma_1,\ldots,\gamma_n\rangle\) with no
expansion marks: it lists each member of \(\Gamma\) exactly once.
\begin{enumerate}
  \item \emph{Satisfiability} of \(\Gamma\) is tested by a tableau rooted at
        \(\Gamma\).
  \item \emph{Consequence} \(\Gamma\entails\varphi\) is tested by a tableau
        rooted at \(\Gamma\cup\{\neg\varphi\}\).  The added formula is
        \(\neg\varphi\) itself, with no simplification even when
        \(\varphi\) is a negation; if \(\neg\varphi\in\Gamma\) already, it is
        listed once.
  \item \emph{Validity} \(\entails\varphi\) is the case
        \(\Gamma=\varnothing\): it is tested by a tableau rooted at
        \(\{\neg\varphi\}\).
\end{enumerate}
\end{hbconvention}

By unique readability (\cref{thm:prop-unique-readability}) every formula of
\(\mathsf{Fm}_{\mathsf{CPL}}\) is a variable \(p\), \(\bottom\), or of the
form \(\neg\chi\) or \((\chi\circ\psi)\) with
\(\circ\in\{\land,\lor,\to\}\), and \(\chi\) in \(\neg\chi\) is again of one of
these forms.  This fixes the following split, on which saturation depends.

\begin{hbdefinition}{def:p02-w02-p02-014-expandable}{Expandable formulas, branches, expansion steps and finished tableaux}
\begin{enumerate}
  \item A formula is \emph{expandable} if it has one of the seven premise
        shapes of the sheet:
        \[
          \varphi\land\psi,\quad \neg(\varphi\land\psi),\quad
          \varphi\lor\psi,\quad \neg(\varphi\lor\psi),\quad
          \varphi\to\psi,\quad \neg(\varphi\to\psi),\quad
          \neg\neg\varphi .
        \]
        Every other formula is \emph{terminal}.  The terminal formulas are
        exactly \(p\), \(\neg p\) (for a variable \(p\)), \(\bottom\) and
        \(\neg\bottom\); the last is the defined verum \(\topv\).  No rule
        has a terminal premise.
  \item Each expandable formula has exactly one rule on the sheet, and that
        rule lists its \emph{alternatives}: one alternative for the
        \emph{linear} shapes and two for the \emph{branching} shapes.
        \[
        \begin{array}{lll}
          \text{premise} & \text{alternatives} & \text{kind}\\ \hline
          \varphi\land\psi & \langle\varphi,\psi\rangle & \text{linear}\\
          \neg(\varphi\lor\psi) & \langle\neg\varphi,\neg\psi\rangle & \text{linear}\\
          \neg(\varphi\to\psi) & \langle\varphi,\neg\psi\rangle & \text{linear}\\
          \neg\neg\varphi & \langle\varphi\rangle & \text{linear}\\
          \neg(\varphi\land\psi) & \langle\neg\varphi\rangle\mid\langle\neg\psi\rangle & \text{branching}\\
          \varphi\lor\psi & \langle\varphi\rangle\mid\langle\psi\rangle & \text{branching}\\
          \varphi\to\psi & \langle\neg\varphi\rangle\mid\langle\psi\rangle & \text{branching}
        \end{array}
        \]
  \item A \emph{branch} is a pair \(B=(\langle\theta_1,\ldots,\theta_m\rangle,M)\)
        of a finite list of formulas and a set \(M\subseteq\{1,\ldots,m\}\)
        of \emph{marked positions}, each holding an expandable formula.  The
        \emph{occurrence} at position \(i\) is the pair \((i,\theta_i)\): two
        positions holding the same formula are two occurrences.  Write
        \(\operatorname{fm}(B)=\{\theta_1,\ldots,\theta_m\}\).  The branch is \emph{closed} if
        \(\bottom\in\operatorname{fm}(B)\) or both \(\chi\in\operatorname{fm}(B)\) and \(\neg\chi\in\operatorname{fm}(B)\) for
        some \(\chi\), and \emph{open} otherwise.
  \item An occurrence at position \(i\) is \emph{expandable on \(B\)} if
        \(B\) is open, \(\theta_i\) is expandable and \(i\notin M\).  An
        \emph{expansion step} at such an occurrence replaces \(B\) by one
        \emph{child} for each alternative \(A\) of \(\theta_i\): the child is
        \((\langle\theta_1,\ldots,\theta_m\rangle\frown A,\;M\cup\{i\})\).
        Nothing is deleted, the new positions are unmarked, and a closed
        branch admits no step.
  \item A \emph{tableau rooted at \(\Gamma\)} is the last member \(T_k\) of an
        \emph{expansion sequence} \(T_0,\ldots,T_k\) of finite lists of
        branches, where \(T_0=\langle B_\Gamma\rangle\) and \(T_{j+1}\)
        replaces one branch of \(T_j\) by its children under one expansion
        step.  It is \emph{closed} if every branch is closed.
  \item An open branch is \emph{saturated} if every position holding an
        expandable formula is marked; terminal occurrences never need a mark.
        A tableau is \emph{finished} if each of its branches is closed or
        saturated.
\end{enumerate}
\end{hbdefinition}

Children keep the positions of their parent, so a position below a split
denotes the same occurrence on every branch descending from the branch where
it was added.  The sheet's instruction to append the conclusions ``on every
open branch through'' an occurrence is therefore a finite run of expansion
steps, one for each open branch that carries that occurrence unmarked, and
every result below applies to it.  The tree drawn in
\cref{fig:p02-closed-tableau} is the picture of an expansion sequence in
which shared initial segments of branches are drawn once.  Saturation is
determinate: whether a branch is saturated depends only on its list and its
mark set.  In particular \(\neg\bottom\) and \(\neg p\) never prevent
saturation, and two occurrences of one expandable formula must each be
expanded.

The termination argument uses a size measure.  For a formula \(\chi\), let
\(\operatorname{sz}(\chi)\) be the number of occurrences of variables,
\(\bottom\) and connectives in \(\chi\) (parentheses are not counted):
\(\operatorname{sz}(p)=\operatorname{sz}(\bottom)=1\),
\(\operatorname{sz}(\neg\chi)=\operatorname{sz}(\chi)+1\) and
\(\operatorname{sz}(\chi\circ\psi)=\operatorname{sz}(\chi)+\operatorname{sz}(\psi)+1\).
For a branch \(B\), put
\[
  \mu(B)\defeq\sum_{\substack{1\leq i\leq m,\ i\notin M,\\
     \theta_i\text{ expandable}}}3^{\operatorname{sz}(\theta_i)},
\]
a sum over positions, so a formula at two unmarked positions counts twice.
For an open branch, \(\mu(B)=0\) exactly
when \(B\) is saturated.  Finally the \emph{branch valuation} of a branch
\(B\) is the valuation \(v_B\) with \(v_B(p)\) true exactly when the variable
\(p\) belongs to \(\operatorname{fm}(B)\).

\subsection{Soundness, termination and the truth lemma}
\label{sec:p02-w02-p02-014-results}

\begin{hblemma}{thm:p02-w02-p02-014-alternatives}{proved=annex:p02-w02-p02-014-alternatives}{Each rule is a semantic case split}
For every expandable formula \(\chi\) and every valuation \(v\),
\(v\entails\chi\) if and only if there is an alternative \(A\) of \(\chi\)
with \(v\entails\alpha\) for every formula \(\alpha\) of \(A\).  Moreover
every formula of every alternative of \(\chi\) has size at most
\(\operatorname{sz}(\chi)-1\), and each alternative has at most two members.
\end{hblemma}

\begin{hblemma}{thm:p02-w02-p02-014-invariants}{proved=annex:p02-w02-p02-014-invariants}{Branch invariants}
Let \(B=(\langle\theta_1,\ldots,\theta_m\rangle,M)\) be a branch of a tableau
rooted at \(\Gamma=\{\gamma_1,\ldots,\gamma_n\}\).
\begin{enumerate}
  \item \(m\geq n\) and \(\theta_\ell=\gamma_\ell\) for \(\ell\leq n\); hence
        \(\Gamma\subseteq\operatorname{fm}(B)\).
  \item For each \(i\in M\), the formula \(\theta_i\) is expandable and some
        alternative \(A\) of \(\theta_i\) satisfies \(A\subseteq\operatorname{fm}(B)\).
\end{enumerate}
\end{hblemma}

\begin{hbtheorem}{thm:p02-w02-p02-014-soundness}{proved=annex:p02-w02-p02-014-soundness}{Soundness of closure}
Let \(\Gamma\) be finite.
\begin{enumerate}
  \item If \(v\entails\operatorname{fm}(B)\) and an expansion step replaces \(B\) by its
        children, then \(v\entails\operatorname{fm}(C)\) for at least one child \(C\).
  \item If \(v\entails\Gamma\), then every tableau rooted at \(\Gamma\) has a
        branch \(B\) with \(v\entails\operatorname{fm}(B)\), and every such branch is open.
  \item If some tableau rooted at \(\Gamma\), finished or not, is closed, then
        \(\Gamma\) is unsatisfiable.
\end{enumerate}
\end{hbtheorem}

\begin{hbtheorem}{thm:p02-w02-p02-014-termination}{proved=annex:p02-w02-p02-014-termination}{Termination of exhaustive expansion}
Let \(\Gamma\) be finite and \(\mu_0\defeq\mu(B_\Gamma)\).
\begin{enumerate}
  \item Every child \(C\) of an expansion step at a branch \(B\) satisfies
        \(\mu(C)<\mu(B)\).
  \item Every expansion sequence \(T_0,\ldots,T_k\) from \(B_\Gamma\) has
        \(k\leq 2^{\mu_0+1}-2\).
  \item A tableau admits an expansion step if and only if it is not finished.
  \item Consequently, expanding expandable occurrences in any order stops
        after at most \(2^{\mu_0+1}-2\) steps, and it stops exactly at a
        finished tableau.  In particular a finished tableau rooted at
        \(\Gamma\) exists.
\end{enumerate}
\end{hbtheorem}

The bound in clause 2 is a crude ceiling that proves termination; it is not
an estimate of the size of a typical tableau.

\begin{hblemma}{thm:p02-w02-p02-014-truth}{proved=annex:p02-w02-p02-014-truth}{Truth lemma for open saturated branches}
Let \(B\) be an open saturated branch of a tableau rooted at a finite set.
Then \(v_B\entails\chi\) for every \(\chi\in\operatorname{fm}(B)\).
\end{hblemma}

\subsection{Completeness, countermodel read-off and decision}
\label{sec:p02-w02-p02-014-completeness}

\begin{hbtheorem}{thm:p02-w02-p02-014-completeness}{proved=annex:p02-w02-p02-014-completeness}{Systematic completeness and countermodel read-off}
Let \(T\) be a finished tableau rooted at a finite set \(\Gamma\).
\begin{enumerate}
  \item \(T\) is closed if and only if \(\Gamma\) is unsatisfiable.
  \item If \(B\) is an open branch of \(T\), then \(v_B\) satisfies every
        formula on \(B\), and in particular \(v_B\entails\Gamma\).
  \item Any two finished tableaux rooted at \(\Gamma\) are either both closed
        or both open, whatever order of expansion produced them.
\end{enumerate}
\end{hbtheorem}

\begin{hbcorollary}{thm:p02-w02-p02-014-consequence}{proved=annex:p02-w02-p02-014-consequence}{Tableau tests for consequence and validity}
Let \(\Gamma\) be finite and \(\varphi\in\mathsf{Fm}_{\mathsf{CPL}}\).
\begin{enumerate}
  \item \(\Gamma\entails\varphi\) if and only if some, equivalently every,
        finished tableau rooted at \(\Gamma\cup\{\neg\varphi\}\) is closed.
  \item If a finished tableau rooted at \(\Gamma\cup\{\neg\varphi\}\) has an
        open branch \(B\), then \(v_B\) is a countervaluation:
        \(v_B\entails\Gamma\) and \(\widehat{v_B}(\varphi)\) is false.
  \item \(\entails\varphi\) if and only if some, equivalently every,
        finished tableau rooted at \(\{\neg\varphi\}\) is closed.
\end{enumerate}
\end{hbcorollary}

\begin{hbmethod}{def:p02-w02-p02-014-procedure}{Systematic tableau procedure}
\emph{Input:} a finite \(\Gamma\subseteq\mathsf{Fm}_{\mathsf{CPL}}\).
\begin{enumerate}
  \item Start with \(T_0=\langle B_\Gamma\rangle\).
  \item While some branch is open and not saturated, take the first such
        branch and the least unmarked position on it holding an expandable
        formula, and perform the expansion step there.
  \item When no step applies, output \emph{unsatisfiable} if every branch is
        closed.  Otherwise take the first open branch \(B\) and output the
        finite set of variables in \(\operatorname{fm}(B)\); it describes \(v_B\).
\end{enumerate}
For \(\Gamma\entails\varphi\) run the procedure on
\(\Gamma\cup\{\neg\varphi\}\) and for \(\entails\varphi\) on
\(\{\neg\varphi\}\): \emph{unsatisfiable} means the consequence (or validity)
holds, and an output valuation is a countervaluation.  Any other fixed rule
for choosing the branch and the position gives the same verdict, by
\cref{thm:p02-w02-p02-014-completeness}(3).
\end{hbmethod}

\begin{hbtheorem}{thm:p02-w02-p02-014-decision}{proved=annex:p02-w02-p02-014-decision}{Propositional tableaux decide by terminating saturation}
For every finite \(\Gamma\), the procedure of
\cref{def:p02-w02-p02-014-procedure} halts after at most
\(2^{\mu(B_\Gamma)+1}-2\) expansion steps.  It outputs \emph{unsatisfiable}
exactly when \(\Gamma\) is unsatisfiable, and any valuation it outputs
satisfies \(\Gamma\).  Hence satisfiability of finite sets of formulas,
validity of formulas, and consequence from finitely many premises are each
decided by an algorithm that halts on every input with the correct answer,
and every negative answer to validity or consequence comes with a
countervaluation.
\end{hbtheorem}

\begin{hbexample}{def:p02-w02-p02-014-example-peirce}{A closed tableau: Peirce's law}
Test \(\entails((P\to Q)\to P)\to P\).  The root branch is
\(\langle\neg(((P\to Q)\to P)\to P)\rangle\).
\begin{enumerate}
  \item Position 1 has the linear shape \(\neg(\varphi\to\psi)\).  Mark it
        and append \((P\to Q)\to P\) (position 2) and \(\neg P\)
        (position 3).
  \item Position 2 has the branching shape \(\varphi\to\psi\).  Mark it.
        The left child appends \(\neg(P\to Q)\) at position 4; the right
        child appends \(P\) at position 4.
  \item The right child contains \(P\) and \(\neg P\) and is closed.
  \item On the left child, position 4 has the shape \(\neg(\varphi\to\psi)\).
        Mark it and append \(P\) and \(\neg Q\).  This child now contains
        \(P\) and \(\neg P\) and is closed.
\end{enumerate}
Both branches are closed, so by \cref{thm:p02-w02-p02-014-consequence} the
formula is valid.  The root measure is
\(\mu_0=3^{\operatorname{sz}(((P\to Q)\to P)\to P)+1}=3^{8}\), and three
steps were used.
\end{hbexample}

\begin{hbexample}{def:p02-w02-p02-014-example-counter}{An open branch and its countervaluation}
Test \(\{P\to Q\}\entails Q\to P\).  The root branch is
\(\langle P\to Q,\ \neg(Q\to P)\rangle\).
\begin{enumerate}
  \item Expand position 2, of shape \(\neg(\varphi\to\psi)\): mark it and
        append \(Q\) (position 3) and \(\neg P\) (position 4).
  \item Expand position 1, of shape \(\varphi\to\psi\): mark it.  The left
        child appends \(\neg P\), the right child appends \(Q\).
\end{enumerate}
The left branch is
\(\langle P\to Q,\neg(Q\to P),Q,\neg P,\neg P\rangle\) with positions 1 and
2 marked.  It is open: \(\bottom\) is not on it, and no formula on it has
its negation on it.  It is saturated, since its unmarked formulas \(Q\) and
\(\neg P\) are terminal.  Its branch valuation makes \(Q\) true and every other variable,
in particular \(P\), false.  It satisfies \(P\to Q\) and falsifies
\(Q\to P\), as \cref{thm:p02-w02-p02-014-consequence}(2) requires, so the
consequence fails.  The right branch is open and saturated as well and gives
the same valuation.  In general different open branches may give different
countervaluations.
\end{hbexample}

\begin{hbwarning}{def:p02-w02-p02-014-unfinished}{An open branch that is not saturated proves nothing}
The truth lemma needs saturation.  The root branch \(\langle P\land\neg
P\rangle\) is open before any expansion, and its branch valuation makes every
variable false; that valuation falsifies \(P\land\neg P\).  Only after the
\(\land\) step does the branch contain \(P\) and \(\neg P\) and close
(\cref{fig:p02-closed-tableau}).  An open branch in an unfinished tableau
therefore supports no conclusion about satisfiability, and a closed branch
alone does not close the tableau: every branch must close.
\end{hbwarning}

\subsection{Why termination is special to the propositional case}
\label{sec:p02-w02-p02-014-fol-contrast}

Termination rests on two facts: an expandable occurrence is expanded at most
once on a branch, and each alternative is smaller than its premise.  The
first-order calculus of \cref{def:p02-fol-tableau} keeps the first fact
for propositional and witness-producing occurrences (its rules are on the
calculus sheet \texttt{sheet:\allowbreak tableaux-fol}), but a universal occurrence, or a negated
existential, stays available at every closed term of the parameter-extended
language \(\mathcal L_{\Gamma}^{*}\) not yet used for it.

\begin{hblemma}{thm:p02-w02-p02-014-fol-no-finite-saturation}{proved=annex:p02-w02-p02-014-fol-no-finite-saturation}{Finite first-order branches with a universal occurrence are not saturated}
Let \(\Gamma\) be a finite set of equality-free sentences and let \(B\) be a
finite open branch of a first-order tableau rooted at \(\Gamma\) in the sense
of \cref{def:p02-fol-tableau}.  If \(B\) contains an occurrence of a formula
\(\forall x\,\varphi\) or \(\neg\exists x\,\varphi\), then \(B\) is not
saturated.  Hence every saturated open branch containing such an occurrence
is infinite.
\end{hblemma}

The rest of this section supplies the first-order metatheory: an exact
bookkeeping for first-order branches, the equality rules, soundness, the
Herbrand truth lemma for open saturated (fair) branches, and the fair
systematic expansion procedure, which semidecides unsatisfiability.

\subsection{First-order tableaux with and without equality}
\label{sec:p02-w02-p02-014-fol-calculus}

Signatures, terms, sentences, structures, satisfaction
\(\mathcal M,s\entails\varphi\) and truth \(\mathcal M\entails\sigma\) are
those of \cref{sec:fol-syntax-semantics-01,sec:fol-semantics-01,sec:fol-semantics-02};
the equality sign is logical and every structure interprets it as identity
on its domain \(\dom\mathcal M\).  For a closed term \(t\), write
\(t^{\mathcal M}\) for its value \(\sem{t}^{\mathcal M}_s\), which does not
depend on the assignment \(s\).  Substitution \(\varphi[t/x]\) is the total
capture-avoiding operation of \cref{sec:fol-syntax-semantics-04}, and
\(c(\varphi)\) is first-order formula complexity, which counts connectives
and quantifiers and is unchanged by substitution
(\cref{lem:substitution-complexity}).

\begin{hbconvention}{def:p02-w02-p02-014-fol-setting}{Root language, parameters and the two calculi}
Let \(\Gamma\) be a finite set of \(\mathcal L\)-sentences.  As in
\cref{def:p02-fol-tableau}, \(\mathcal L_{\Gamma}\) contains exactly the
function and relation symbols occurring in \(\Gamma\), and
\(\mathcal L_{\Gamma}^{*}\) adds pairwise distinct new constant symbols
\(a_0,a_1,\ldots\), the \emph{parameters}, none of which belongs to
\(\mathcal L\).  Let \(\mathrm{CT}_{\Gamma}\) be the set of closed
\(\mathcal L_{\Gamma}^{*}\)-terms.  It is countably infinite, since
\(\mathcal L_{\Gamma}\) is finite and the parameters are infinitely many;
fix a computable enumeration \(t_0,t_1,t_2,\ldots\) in which every member of
\(\mathrm{CT}_{\Gamma}\) occurs.
\begin{enumerate}
  \item The \emph{equality-free calculus} \(\mathsf T\) is the calculus of
        \cref{def:p02-fol-tableau}, whose rules are on the calculus sheet
        \texttt{sheet:\allowbreak tableaux-fol}.  It is used only when no
        member of \(\Gamma\) contains the equality sign.
  \item The \emph{calculus with equality} \(\mathsf T^{=}\) has the rules
        of \(\mathsf T\) together with the two equality rules of
        \cref{def:p02-w02-p02-014-fol-equality}.  Its root \(\Gamma\) may
        contain equations.
\end{enumerate}
By unique readability (\cref{thm:fol-unique-readability}) every
\(\mathcal L_{\Gamma}^{*}\)-sentence has exactly one of the following
\emph{kinds}.  It is \emph{propositional-expandable} if it has one of the
seven premise shapes of \cref{def:p02-w02-p02-014-expandable}, with the
alternatives listed there; a \emph{gamma} sentence if it is
\(\forall x\,\varphi\) or \(\neg\exists x\,\varphi\); a \emph{delta}
sentence if it is \(\exists x\,\varphi\) or \(\neg\forall x\,\varphi\); and
\emph{terminal} otherwise.  The terminal sentences are exactly the atomic
sentences \(R(t_1,\ldots,t_n)\) and \(r=u\), their negations,
\(\bottom\) and \(\neg\bottom\).
\end{hbconvention}

\begin{hbdefinition}{def:p02-w02-p02-014-fol-equality}{Tableau equality rules}
The calculus \(\mathsf T^{=}\) adds two linear rules, with closed terms
\(r,u,t\in\mathrm{CT}_{\Gamma}\).
\begin{enumerate}
  \item \emph{Reflexivity.}  For any \(t\in\mathrm{CT}_{\Gamma}\) not yet
        used for this rule on the branch, append \(t=t\).  The rule has no
        premise; like a gamma rule it stays available at every unused term.
  \item \emph{Replacement.}  Let an occurrence of the equation \(r=u\) and an
        occurrence of an atomic sentence \(\alpha\) lie on the branch, and let
        \(\varphi\) be an atomic \(\mathcal L_{\Gamma}^{*}\)-formula in which
        the variable \(x_0\) occurs exactly once and no other variable
        occurs, with \(\varphi[r/x_0]=\alpha\).  If this triple has not yet
        been used on the branch, append \(\varphi[u/x_0]\).
\end{enumerate}
Replacement therefore rewrites one occurrence of \(r\) in \(\alpha\) to
\(u\), left to right; \(\alpha\) may itself be an equation, and the two
occurrences may coincide.  For given occurrences of \(r=u\) and
\(\alpha\) there are as many choices of \(\varphi\) as there are
occurrences of \(r\) as a subterm of \(\alpha\), which is finitely many.
The rules are the tableau forms of the reflexivity and replacement rules of
the equality calculus of \texttt{sheet:\allowbreak nd-fol}
(\cref{thm:p02-equality-laws}); symmetry and transitivity are not rules here
but are recovered from these two on every saturated branch
(\cref{thm:p02-w02-p02-014-congruence}).
\end{hbdefinition}

\begin{hbdefinition}{def:p02-w02-p02-014-fol-branch}{First-order branches, tasks and steps}
Fix \(\Gamma\) and one of the calculi \(\mathsf T\), \(\mathsf T^{=}\).
\begin{enumerate}
  \item A \emph{branch} is a pair
        \(B=(\langle\theta_1,\theta_2,\ldots\rangle,D)\) of a finite or
        infinite list of \(\mathcal L_{\Gamma}^{*}\)-sentences and a set
        \(D\) of \emph{done tasks}.  Write \(\operatorname{fm}(B)\) for the
        set of sentences on the list.  Closed and open are as in
        \cref{def:p02-w02-p02-014-expandable}: \(B\) is \emph{closed} if
        \(\bottom\in\operatorname{fm}(B)\) or \(\chi,\neg\chi\in\operatorname{fm}(B)\) for some \(\chi\).
  \item The \emph{tasks available on \(B\)} are the following.
        \begin{itemize}
          \item \(\langle i\rangle\), for each position \(i\) holding a
                delta sentence or a propositional-expandable one.
          \item \(\langle i,t\rangle\), for each position \(i\) holding a
                gamma sentence and each \(t\in\mathrm{CT}_{\Gamma}\).
          \item In \(\mathsf T^{=}\) only: \(\langle{=},t\rangle\) for each
                \(t\in\mathrm{CT}_{\Gamma}\); and
                \(\langle i,j,\varphi\rangle\) for positions \(i,j\) such
                that \(\theta_i\) is an equation \(r=u\) and \(\varphi\) is
                as in the replacement rule with
                \(\varphi[r/x_0]=\theta_j\).
        \end{itemize}
  \item An \emph{expansion step} applies to a finite open branch
        \(B=(\langle\theta_1,\ldots,\theta_m\rangle,D)\) at an available task
        \(\tau\notin D\).  It replaces \(B\) by one child
        \((\langle\theta_1,\ldots,\theta_m\rangle\frown A,\;D\cup\{\tau\})\)
        for each alternative \(A\) determined as follows.  For
        \(\langle i\rangle\) with \(\theta_i\) propositional-expandable, the
        alternatives of \(\theta_i\).  For \(\langle i\rangle\) with
        \(\theta_i=\exists x\,\varphi\) or \(\neg\forall x\,\varphi\), the
        single alternative \(\langle\varphi[a/x]\rangle\) or
        \(\langle\neg(\varphi[a/x])\rangle\), where \(a\) is the parameter of
        least index not occurring in \(\theta_1,\ldots,\theta_m\).  For
        \(\langle i,t\rangle\) with \(\theta_i=\forall x\,\varphi\) or
        \(\neg\exists x\,\varphi\), the single alternative
        \(\langle\varphi[t/x]\rangle\) or \(\langle\neg(\varphi[t/x])\rangle\).
        For \(\langle{=},t\rangle\), \(\langle t=t\rangle\); for
        \(\langle i,j,\varphi\rangle\) with \(\theta_i\) the equation \(r=u\),
        \(\langle\varphi[u/x_0]\rangle\).
  \item The \emph{root branch} \(B_\Gamma\) lists the members of \(\Gamma\) once
        each, in a fixed order, with \(D=\varnothing\).  An \emph{expansion
        sequence} is a finite or infinite sequence \(T_0,T_1,\ldots\) of
        finite lists of finite branches with \(T_0=\langle B_\Gamma\rangle\),
        each \(T_{j+1}\) replacing one branch of \(T_j\) by its children under
        one step.  Each branch of \(T_{j+1}\) has a \emph{parent} in \(T_j\):
        the replaced branch, if it is one of its children, and itself
        otherwise.  Iterating parents gives, for \(j\leq j'\), the unique
        \emph{ancestor} in \(T_j\) of each branch of \(T_{j'}\); a branch
        \emph{descends} from its ancestors.  A \emph{finite tableau} rooted at
        \(\Gamma\) is the last member of a finite expansion sequence; it is
        \emph{closed} if all its branches are closed.
  \item A \emph{path} through an infinite expansion sequence is a sequence
        \((B_j)_{j\in\N}\) in which \(B_j\) is a branch of \(T_j\) and
        \(B_j\) is the parent of \(B_{j+1}\).  Its \emph{limit branch}
        \(B_\omega\) has as list the union of the lists of the \(B_j\), each
        of which extends the previous one, and as done set the union of
        their done sets.  The finite branches of the \(T_j\) and the limit
        branches are the branches of the tableau, possibly infinite,
        generated by the sequence.
  \item An open branch \(B\), finite or limit, is \emph{saturated} if every
        task available on \(B\) belongs to \(D\).
\end{enumerate}
\end{hbdefinition}

The tasks make the sheet's records exact.  A done \(\langle i\rangle\) is the
mark of a propositional or delta occurrence.  The done tasks
\(\langle i,t\rangle\) record the terms used for the gamma occurrence at
position \(i\), so a saturated branch has used every \(t_k\) for every gamma
occurrence: saturation in the sense of item 6 is exactly ``fair and
saturated'' in the sense of \cref{def:p02-fol-tableau}.  The freshness
condition is the choice of the least unused parameter in a delta step, and
the instantiation condition is that gamma and reflexivity steps range over
all of \(\mathrm{CT}_{\Gamma}\), not only over terms already on the branch.
As in the propositional case, the sheet's instruction to append conclusions
``on every open branch through'' an occurrence is a finite run of single
steps.  Every instance \(\varphi[t/x]\) appended by a step is again a
sentence, because the premise is a sentence and \(t\) is closed
(\cref{lem:substitution-fv}).  In \(\mathsf T^{=}\) no finite branch is
saturated, since infinitely many reflexivity tasks are available on it.

\begin{hblemma}{thm:p02-w02-p02-014-fol-cases}{proved=annex:p02-w02-p02-014-fol-cases}{First-order rules are sound case splits}
Let \(\mathcal M\) be an \(\mathcal L_{\Gamma}^{*}\)-structure.
\begin{enumerate}
  \item For every propositional-expandable sentence \(\chi\),
        \(\mathcal M\entails\chi\) iff \(\mathcal M\) satisfies every member
        of some alternative of \(\chi\).
  \item If \(\mathcal M\entails\forall x\,\varphi\), then
        \(\mathcal M\entails\varphi[t/x]\) for every
        \(t\in\mathrm{CT}_{\Gamma}\); if
        \(\mathcal M\entails\neg\exists x\,\varphi\), then
        \(\mathcal M\entails\neg(\varphi[t/x])\) for every such \(t\).
  \item Let \(a\) be a parameter that occurs in no member of a set \(X\) of
        sentences.  If \(\mathcal M\entails X\) and
        \(\exists x\,\varphi\in X\), there is a structure \(\mathcal M'\),
        differing from \(\mathcal M\) at most in the interpretation of
        \(a\), with \(\mathcal M'\entails X\cup\{\varphi[a/x]\}\).  The same
        holds with \(\neg\forall x\,\varphi\in X\) and
        \(\neg(\varphi[a/x])\).
  \item \(\mathcal M\entails t=t\); and if \(\mathcal M\entails r=u\) and
        \(\mathcal M\entails\varphi[r/x_0]\) for \(\varphi\) as in the
        replacement rule, then \(\mathcal M\entails\varphi[u/x_0]\).
\end{enumerate}
\end{hblemma}

\begin{hblemma}{thm:p02-w02-p02-014-fol-invariants}{proved=annex:p02-w02-p02-014-fol-invariants}{First-order branch invariants}
Let \(B\) be a finite or limit branch generated from \(B_\Gamma\).  Then
\(\Gamma\subseteq\operatorname{fm}(B)\), and every done task has its conclusion on \(B\):
for \(\langle i\rangle\) with \(\theta_i\) propositional-expandable, some
alternative of \(\theta_i\) is included in \(\operatorname{fm}(B)\); for
\(\langle i\rangle\) with \(\theta_i=\exists x\,\varphi\) (or
\(\neg\forall x\,\varphi\)), \(\varphi[a/x]\in\operatorname{fm}(B)\) (or
\(\neg(\varphi[a/x])\in\operatorname{fm}(B)\)) for some parameter \(a\); for
\(\langle i,t\rangle\), the instance of \(\theta_i\) at \(t\) is in \(\operatorname{fm}(B)\);
for \(\langle{=},t\rangle\), \(t=t\in\operatorname{fm}(B)\); and for
\(\langle i,j,\varphi\rangle\) with \(\theta_i\) the equation \(r=u\),
\(\varphi[u/x_0]\in\operatorname{fm}(B)\).  In \(\mathsf T\), no member of \(\operatorname{fm}(B)\)
contains the equality sign.  A limit branch is open if and only if every
\(B_j\) on its path is open.
\end{hblemma}

\begin{hbtheorem}{thm:p02-w02-p02-014-fol-soundness}{proved=annex:p02-w02-p02-014-fol-soundness}{First-order soundness of closure}
Let \(\Gamma\) be a finite set of \(\mathcal L\)-sentences, equality-free
for \(\mathsf T\).
\begin{enumerate}
  \item If an \(\mathcal L_{\Gamma}^{*}\)-structure satisfies \(\operatorname{fm}(B)\) and
        an expansion step replaces the finite branch \(B\), then some
        \(\mathcal L_{\Gamma}^{*}\)-structure satisfies \(\operatorname{fm}(C)\) for some
        child \(C\).
  \item If \(\Gamma\) has a model, every finite tableau rooted at \(\Gamma\)
        has an open branch.
  \item If some finite tableau rooted at \(\Gamma\) is closed, then
        \(\Gamma\) is unsatisfiable.
\end{enumerate}
\end{hbtheorem}

\subsection{Hintikka sets and the Herbrand truth lemma}
\label{sec:p02-w02-p02-014-fol-herbrand}

\begin{hbdefinition}{def:p02-w02-p02-014-hintikka}{Hintikka sets and their quotient Herbrand structures}
A set \(H\) of \(\mathcal L_{\Gamma}^{*}\)-sentences is a \emph{Hintikka
set} if:
\begin{enumerate}
  \item[(H0)] \(\bottom\notin H\), and no \(\chi\) has both \(\chi\in H\)
        and \(\neg\chi\in H\);
  \item[(H1)] if \(\chi\in H\) is propositional-expandable, every member of
        some alternative of \(\chi\) is in \(H\);
  \item[(H2)] if \(\forall x\,\varphi\in H\), then \(\varphi[t/x]\in H\), and
        if \(\neg\exists x\,\varphi\in H\), then \(\neg(\varphi[t/x])\in H\),
        for every \(t\in\mathrm{CT}_{\Gamma}\);
  \item[(H3)] if \(\exists x\,\varphi\in H\), then \(\varphi[t/x]\in H\), and
        if \(\neg\forall x\,\varphi\in H\), then \(\neg(\varphi[t/x])\in H\),
        for some \(t\in\mathrm{CT}_{\Gamma}\);
  \item[(H4)] if some member of \(H\) contains the equality sign, then
        \(t=t\in H\) for every \(t\in\mathrm{CT}_{\Gamma}\), and
        \(\varphi[u/x_0]\in H\) whenever \(r=u\in H\) and
        \(\varphi[r/x_0]\in H\) for \(\varphi\) as in the replacement rule.
\end{enumerate}
\(H\) is \emph{equality-free} if none of its members contains the equality
sign.  For a Hintikka set \(H\), put \(r\sim_H u\) iff \(r\) and \(u\) are
the same term or \(r=u\in H\).  The \emph{quotient Herbrand structure}
\(\mathcal M_H\) has domain \(\mathrm{CT}_{\Gamma}/{\sim_H}\), whose
members are the classes \([t]\); it interprets \(f\) by
\(f^{\mathcal M_H}([t_1],\ldots,[t_n])=[f(t_1,\ldots,t_n)]\) and \(R\) by
\(R^{\mathcal M_H}=\{([t_1],\ldots,[t_n]):R(t_1,\ldots,t_n)\in H\}\).
\end{hbdefinition}

\begin{hblemma}{thm:p02-w02-p02-014-congruence}{proved=annex:p02-w02-p02-014-congruence}{The branch equality is a congruence}
Let \(H\) be a Hintikka set.  Then \(\sim_H\) is an equivalence relation on
\(\mathrm{CT}_{\Gamma}\); if \(t_i\sim_H u_i\) for \(1\leq i\leq n\), then
\(f(t_1,\ldots,t_n)\sim_H f(u_1,\ldots,u_n)\), and
\(R(t_1,\ldots,t_n)\in H\) implies \(R(u_1,\ldots,u_n)\in H\).  Hence
\(\mathcal M_H\) is well defined, \(t^{\mathcal M_H}=[t]\) for every
\(t\in\mathrm{CT}_{\Gamma}\), and
\(([t_1],\ldots,[t_n])\in R^{\mathcal M_H}\) iff
\(R(t_1,\ldots,t_n)\in H\).  If \(H\) is equality-free, \(\sim_H\) is
identity, each class is a singleton, and \(\mathcal M_H\) is, after
identifying \([t]\) with \(t\), the Herbrand structure
\(\mathcal M_{B^{+}_{\mathrm{at}}}\) of
\cref{thm:p02-fol-tableau-adequacy} with \(B^{+}_{\mathrm{at}}\) the atomic
members of \(H\).
\end{hblemma}

\begin{hblemma}{thm:p02-w02-p02-014-herbrand-truth}{proved=annex:p02-w02-p02-014-herbrand-truth}{Herbrand truth lemma}
If \(H\) is a Hintikka set, then \(\mathcal M_H\entails\chi\) for every
\(\chi\in H\).
\end{hblemma}

\begin{hblemma}{thm:p02-w02-p02-014-saturated-hintikka}{proved=annex:p02-w02-p02-014-saturated-hintikka}{Open saturated branches are Hintikka sets}
If \(B\) is an open saturated branch, finite or limit, of a tableau rooted at
\(\Gamma\) in \(\mathsf T\) or \(\mathsf T^{=}\), then \(\operatorname{fm}(B)\) is a Hintikka
set, and in \(\mathsf T\) it is equality-free.
\end{hblemma}

\begin{hbtheorem}{thm:p02-w02-p02-014-fol-extraction}{proved=annex:p02-w02-p02-014-fol-extraction}{Countermodel extraction from an open saturated branch}
Let \(B\) be an open saturated branch, finite or limit, of a tableau rooted
at a finite set \(\Gamma\) of \(\mathcal L\)-sentences in \(\mathsf T\) or
\(\mathsf T^{=}\), and let \(H=\operatorname{fm}(B)\).  Then
\(\mathcal M_H\entails\chi\) for every \(\chi\) on \(B\).  Its reduct
\(\mathcal M_H\restriction\mathcal L_{\Gamma}\) satisfies \(\Gamma\), and
interpreting the symbols of \(\mathcal L\) outside \(\mathcal L_{\Gamma}\)
arbitrarily on the same domain gives an \(\mathcal L\)-structure satisfying
\(\Gamma\).  In \(\mathsf T\), \(\mathcal M_H\) is the Herbrand structure
\(\mathcal M_{B^{+}_{\mathrm{at}}}\) of
\cref{thm:p02-fol-tableau-adequacy}.
\end{hbtheorem}

\subsection{Fair systematic expansion and semidecision}
\label{sec:p02-w02-p02-014-fol-systematic}

An arbitrary order of expansion can starve a gamma occurrence forever
(\cref{def:p02-w02-p02-014-fol-unfair}).  The following procedure performs
at each stage a finite batch of tasks that grows with the stage, so every
task available on a surviving branch is eventually performed.

\begin{hbmethod}{def:p02-w02-p02-014-fol-procedure}{Fair systematic expansion}
\emph{Input:} a finite set \(\Gamma\) of \(\mathcal L\)-sentences, with the
calculus \(\mathsf T\) if \(\Gamma\) is equality-free and \(\mathsf T^{=}\)
otherwise, and the enumeration \(t_0,t_1,\ldots\) of
\cref{def:p02-w02-p02-014-fol-setting}.  Let \(T^{(0)}=\langle
B_\Gamma\rangle\).  At stage \(n=0,1,2,\ldots\), with current finite tableau
\(T^{(n)}\):
\begin{enumerate}
  \item If every branch of \(T^{(n)}\) is closed, halt with output
        \emph{unsatisfiable} and the closed tableau \(T^{(n)}\).
  \item If some branch of \(T^{(n)}\) is finite, open and saturated, halt
        with output \emph{satisfiable} and that branch.
  \item Otherwise, for each branch \(B\) of \(T^{(n)}\) that is open, with
        list length \(m_B\), form the finite list \(W_n(B)\) of the
        following tasks, in a fixed order: \(\langle i\rangle\) for
        \(i\leq m_B\) holding a propositional-expandable or delta sentence;
        \(\langle i,t_k\rangle\) for \(i\leq m_B\) holding a gamma sentence
        and \(k\leq n\); and, in \(\mathsf T^{=}\), \(\langle{=},t_k\rangle\)
        for \(k\leq n\) and every replacement task \(\langle
        i,j,\varphi\rangle\) with \(i,j\leq m_B\).  Taking the tasks of
        \(W_n(B)\) in order, perform the expansion step at the task on every
        open branch that descends from \(B\) and does not yet have that task
        done, until none remains.  The resulting tableau is
        \(T^{(n+1)}\).
\end{enumerate}
\end{hbmethod}

\begin{hbtheorem}{thm:p02-w02-p02-014-fol-semidecision}{proved=annex:p02-w02-p02-014-fol-semidecision}{Fair systematic expansion semidecides unsatisfiability}
Let \(\Gamma\) be a finite set of \(\mathcal L\)-sentences and run
\cref{def:p02-w02-p02-014-fol-procedure}.
\begin{enumerate}
  \item Each stage performs finitely many expansion steps, and the steps of
        all stages, in order, form an expansion sequence from \(B_\Gamma\).
  \item \emph{Infinite branch construction.}  If the procedure does not
        halt, there is a path \((B_j)_{j\in\N}\) through this expansion
        sequence whose limit branch \(B_\omega\) is open and saturated, and
        so fair; \(B_\omega\) is infinite.
  \item The procedure halts with output \emph{unsatisfiable} if and only if
        \(\Gamma\) is unsatisfiable, and then it outputs a finite closed
        tableau rooted at \(\Gamma\).  It halts with output
        \emph{satisfiable} only if \(\Gamma\) is satisfiable, and when
        \(\Gamma\) is satisfiable and it does not halt, the model
        \(\mathcal M_{\operatorname{fm}(B_\omega)}\) of clause 2 satisfies \(\Gamma\) by
        \cref{thm:p02-w02-p02-014-fol-extraction}.
  \item Consequently the procedure is an algorithm that halts with output
        \emph{unsatisfiable} on exactly the unsatisfiable finite sets of
        sentences.
\end{enumerate}
\end{hbtheorem}

\begin{hbcorollary}{thm:p02-w02-p02-014-fol-completeness}{proved=annex:p02-w02-p02-014-fol-completeness}{First-order tableau completeness, consequence and validity}
Let \(\Gamma\) be a finite set of \(\mathcal L\)-sentences and \(\psi\) an
\(\mathcal L\)-sentence; use \(\mathsf T\) when the sentences involved are
equality-free and \(\mathsf T^{=}\) otherwise.
\begin{enumerate}
  \item \(\Gamma\) is unsatisfiable if and only if some finite closed
        tableau rooted at \(\Gamma\) exists.
  \item \(\Gamma\entails\psi\) if and only if some finite closed tableau
        rooted at \(\Gamma\cup\{\neg\psi\}\) exists, and \(\entails\psi\) if
        and only if some finite closed tableau rooted at \(\{\neg\psi\}\)
        exists.
  \item If an open saturated branch \(B\) of a tableau rooted at
        \(\Gamma\cup\{\neg\psi\}\) exists, the expansion of
        \(\mathcal M_{\operatorname{fm}(B)}\restriction\mathcal L_{\Gamma\cup\{\neg\psi\}}\) of
        \cref{thm:p02-w02-p02-014-fol-extraction} is a countermodel:
        it satisfies \(\Gamma\) and falsifies \(\psi\).
\end{enumerate}
\end{hbcorollary}

\begin{hbexample}{def:p02-w02-p02-014-fol-example-closed}{A closed first-order tableau}
In \(\mathsf T\), test the consequence
\[
  \{\forall x(Px\to Qx),\ \exists x\,Px\}\entails\exists x\,Qx .
\]
The root branch is
\(\langle\forall x(Px\to Qx),\ \exists x\,Px,\ \neg\exists x\,Qx\rangle\);
\(\mathcal L_{\Gamma}\) has no function symbols, so the closed terms are the
parameters.
\begin{enumerate}
  \item Task \(\langle 2\rangle\), a delta step: the least parameter not on
        the branch is \(a_0\); append \(Pa_0\) (position 4).
  \item Task \(\langle 1,a_0\rangle\), a gamma step: append
        \(Pa_0\to Qa_0\) (position 5).
  \item Task \(\langle 3,a_0\rangle\), a gamma step: append \(\neg Qa_0\)
        (position 6).
  \item Task \(\langle 5\rangle\), branching: the left child appends
        \(\neg Pa_0\) and closes against position 4; the right child
        appends \(Qa_0\) and closes against position 6.
\end{enumerate}
The finite tableau is closed, so the consequence holds by
\cref{thm:p02-w02-p02-014-fol-completeness}(2).  The gamma steps were made at
the parameter that the delta step introduced.  The systematic procedure
need not choose this order, but fairness guarantees that it performs both
gamma tasks at \(a_0\) by stage \(k\), where \(a_0\) is \(t_k\) in the
enumeration.
\end{hbexample}

\begin{hbexample}{def:p02-w02-p02-014-fol-example-eq}{Equality steps}
Let \(c,d\) be constant symbols and \(\Gamma=\{c=d,\ Pd,\ \neg Pc\}\), in
\(\mathsf T^{=}\).
\begin{enumerate}
  \item Task \(\langle{=},c\rangle\): append \(c=c\) (position 4).
  \item Task \(\langle 1,4,\varphi\rangle\) with \(\varphi\) the formula
        \(x_0=c\): since \(\varphi[c/x_0]\) is \(c=c\), append
        \(\varphi[d/x_0]\), that is \(d=c\) (position 5).
  \item Task \(\langle 5,2,Px_0\rangle\): since \((Px_0)[d/x_0]\) is \(Pd\),
        append \(Pc\) (position 6).  The branch now contains \(Pc\) and
        \(\neg Pc\) and closes.
\end{enumerate}
Replacement rewrites only from left to right, so symmetry had to be
obtained from reflexivity first, exactly as in the proof of
\cref{thm:p02-w02-p02-014-congruence}.
\end{hbexample}

\begin{hbexample}{def:p02-w02-p02-014-fol-example-open}{A finite saturated open branch and its Herbrand countermodel}
Let \(c\) be a constant symbol and \(\Gamma=\{\exists x\,Px,\ \neg Pc\}\), in
\(\mathsf T\).  At stage 0 the only task is \(\langle 1\rangle\); the least
parameter not on the branch is \(a_0\), and the step appends \(Pa_0\).  At
stage 1 the branch \(\langle\exists x\,Px,\neg Pc,Pa_0\rangle\) is open and,
as it carries no gamma sentence, saturated, so the procedure halts with
output \emph{satisfiable}.  The structure \(\mathcal M_H\) has domain the
closed terms \(c,a_0,a_1,\ldots\) and interprets \(P\) as \(\{a_0\}\); its
reduct to \(\{P,c\}\) satisfies \(\exists x\,Px\) and \(\neg Pc\).
\end{hbexample}

\begin{hbexample}{def:p02-w02-p02-014-fol-example}{A satisfiable first-order root with no finite finished tableau}
Let \(\Gamma=\{\forall x\,Px\}\) with \(P\) a unary relation symbol.  Then
\(\mathcal L_{\Gamma}\) has no constant or function symbol, so the closed
\(\mathcal L_{\Gamma}^{*}\)-terms are exactly the parameters \(a_0,a_1,\ldots\).
The root occurrence is retained on every branch, so by
\cref{thm:p02-w02-p02-014-fol-no-finite-saturation} no finite tableau rooted
at \(\Gamma\) has a saturated open branch, and no finite tableau rooted at
\(\Gamma\) is closed, since \(\Gamma\) is true in a one-element structure
in which \(P\) holds of that element
(\cref{thm:p02-w02-p02-014-fol-soundness}(3)).  The procedure of
\cref{def:p02-w02-p02-014-fol-procedure} therefore never halts.  Stage \(n\)
appends \(Pt_n\); the limit branch
\(\langle\forall x\,Px,Pt_0,Pt_1,\ldots\rangle\) is open, saturated
and infinite, and
\(\mathcal M_H\) has domain \(\{a_0,a_1,\ldots\}\) with \(P\) true of every
element.  So the propositional stopping rule ``expand until every branch is
closed or saturated'' never stops here.
\end{hbexample}

\begin{hbwarning}{def:p02-w02-p02-014-fol-unfair}{An unfair branch is not a countermodel}
Let \(\Gamma=\{\forall x\,Px,\ \forall x\,\neg Px\}\), which is
unsatisfiable.  Expanding only the first occurrence, at
\(a_0,a_1,a_2,\ldots\) in turn, is a legal infinite expansion sequence whose
single limit branch \(\langle\forall x\,Px,\forall x\,\neg
Px,Pa_0,Pa_1,\ldots\rangle\) is open.  It is not saturated: no task
\(\langle 2,t\rangle\) is done.  The Herbrand structure whose relations are given
by the atoms on the branch makes \(P\) true of every parameter and falsifies
\(\forall x\,\neg Px\).  The systematic
procedure performs \(\langle 2,t_0\rangle\) at stage 0 and closes the
tableau.  Openness without fairness supports no conclusion.
\end{hbwarning}

\begin{hbwarning}{def:p02-w02-p02-014-fol-freshness}{The delta parameter must be new}
The set \(\Gamma\) of
\cref{def:p02-w02-p02-014-fol-example-open}, with members \(\exists x\,Px\)
and \(\neg Pc\), is satisfiable.  A delta
step that used \(c\) instead of a new parameter would append \(Pc\) and close
the only branch, a closed tableau for a satisfiable set.  The proof of
\cref{thm:p02-w02-p02-014-fol-cases}(3) uses exactly that the parameter
does not occur on the branch.
\end{hbwarning}

\begin{proofboundary}{def:p02-w02-p02-014-fol-boundary}{Decision versus semidecision}
For a finite propositional input, \cref{thm:p02-w02-p02-014-decision}
gives a procedure that always halts and answers both ways.  For first-order
input, \cref{thm:p02-w02-p02-014-fol-semidecision} gives a procedure that
halts on every unsatisfiable input with a finite closed tableau, and that may
run forever on a satisfiable one, as in
\cref{def:p02-w02-p02-014-fol-example}; on a satisfiable input it builds an
infinite fair branch whose quotient Herbrand structure is a model.  The
difference is not an artefact of the procedure: by
\cref{thm:fol-validity-undecidable}, whose status is cited there, first-order
validity is not decidable, and \(\psi\) is valid exactly when \(\{\neg\psi\}\)
is unsatisfiable, so no algorithm halts on every input with the correct
satisfiability verdict.  This section proves the semidecision, the infinite
branch construction and the failure of the propositional termination
argument (\cref{thm:p02-w02-p02-014-fol-no-finite-saturation}); the
undecidability theorem is used only for this concluding remark.
\end{proofboundary}

\subsection{Relationships and sources}
\label{sec:p02-w02-p02-014-relations}

The tableau verdict agrees with the other classical propositional calculi:
by \cref{thm:p02-w02-p02-014-consequence} and the derivability equivalence
\cref{thm:p02-calculi-derivability-equivalence}, whose proof is cited there,
a finite consequence has a
closed tableau exactly when it is derivable in each of the calculi compared
in \cref{tab:p02-proof-systems}.  Resolution
(\cref{thm:prop-resolution-complete}) is also a refutation method, but it
saturates a clause set instead of splitting branches and does not keep a
branch from which a valuation is read.  A finished tableau for
\(\{\neg\varphi\}\) with an open branch is a semantic certificate that
\(\varphi\) is not valid; for formulas of the \(\{\neg,\to\}\) language, the rejection calculus of
\cref{sec:p02-classical-rejection} certifies the same fact by a syntactic
derivation of the rejection judgment, which is a different kind of object.

In the first-order case the Herbrand truth lemma plays the role of the truth
lemma for maximal consistent sets with witnesses
(\cref{lem:fol-truth}): both build a model from closed terms modulo a
provable or listed equality.  The difference is that a Hintikka set need not
be maximal, and the tableau constructs it by a fair search instead of by an
enumeration of all sentences.  The first-order results of this section prove
every clause of \cref{thm:p02-fol-tableau-adequacy} and extend it to roots
with equality.

The unsigned uniform rule format on the sheets is attributed there to
\cite{Smullyan1968}.  The definitions and proofs in this section, including
the task bookkeeping, the tableau equality rules derived from the
natural-deduction equality rules, and the staged fair procedure, are written
directly for the retained-occurrence formats of the calculus sheets
\texttt{sheet:\allowbreak tableaux-cpl} and \texttt{sheet:\allowbreak
tableaux-fol}; they are not a transcription of a published proof.
